%%%%%%%%%%%%%%%%%%%%%%%%%%
%%% 08-samuel/1-24.tex %%%
%%%%%%%%%%%%%%%%%%%%%%%%%%

\Chapter{24}

\Verse{1}
{
	\hX{וַיַּעַל דָּוִד מִשָּׁם וַיֵּשֶׁב בִּמְצָדוֹת עֵין־גֶּדִי׃}
}
{
	\eX{
		David went from there and stayed in the wildernesses of En-gedi.
	}
}
{
	\aB{
		David goes to a different wilderness.
	}
}

\Verse{2}
{
	\hX{וַיְהִי כַּאֲשֶׁר שָׁב שָׁאוּל מֵאַחֲרֵי פְּלִשְׁתִּים וַיַּגִּדוּ לוֹ לֵאמֹר הִנֵּה דָוִד בְּמִדְבַּר עֵין גֶּדִי׃}
}
{
	\eX{
		When Saul returned from pursuing the Philistines, he was told that
		David was in the wilderness of En-gedi.
	}
}
{
	\aB{
		When Saul gets back, he's told that David is in the wilderness of Spring Of The Young Goat.
	}
}

\Verse{3}
{
	\hX{וַיִּקַּח שָׁאוּל שְׁלֹשֶׁת אֲלָפִים אִישׁ בָּחוּר מִכּל־יִשְׂרָאֵל וַיֵּלֶךְ לְבַקֵּשׁ אֶת־דָּוִד וַאֲנָשָׁיו עַל־פְּנֵי צוּרֵי הַיְּעֵלִים׃}
}
{
	\eX{
		So Saul took three thousand picked men from all Israel and went in
		search of David and his men in the direction of the rocks of the
		wild goats;
	}
}
{
	\aB{
		Saul takes 3,000 guys and looks for David near the goats.
	}
}

\Verse{4}
{
	\hX{וַיָּבֹא אֶל־גִּדְרוֹת הַצֹּאן עַל־הַדֶּרֶךְ וְשָׁם מְעָרָה וַיָּבֹא שָׁאוּל לְהָסֵךְ אֶת־רַגְלָיו וְדָוִד וַאֲנָשָׁיו בְּיַרְכְּתֵי הַמְּעָרָה יֹשְׁבִים׃}
}
{
	\eX{
		and he came to the sheepfolds along the way. There was a cave
		there, and Saul went in to relieve himself. Now David and his men
		were sitting in the back of the cave.
	}
}
{
	\aB{
		Saul comes to a sheep pen with a cave near it.
	}

	\aB{
		Saul goes into the cave to poop.
	}

	\aB{
		David and his men are secretly sitting in the pooping cave already, in the back.
	}
}

\Verse{5}
{
	\hX{
		וַיֹּאמְרוּ אַנְשֵׁי דָוִד אֵלָיו הִנֵּה הַיּוֹם אֲשֶׁר־אָמַר יְהֹוָה אֵלֶיךָ הִנֵּה אָנֹכִי נֹתֵן אֶת־[אֹיִבְךָ] (איביך) בְּיָדֶךָ וְעָשִׂיתָ לּוֹ כַּאֲשֶׁר יִטַב בְּעֵינֶיךָ
		וַיָּקם דָּוִד וַיִּכְרֹת אֶת־כְּנַף־הַמְּעִיל אֲשֶׁר־לְשָׁאוּל בַּלָּט׃
	}
}
{
	\eX{
		David’s men said to him, “This is the day of which YHWH said to
		you, ‘I will deliver your enemy into your hands; you can do with
		him as you please.’” David went and stealthily cut off the corner
		of Saul’s cloak.
	}
}
{
	Qere Ketiv.

	\aB{
		David's men are like ``We got him!''
	}

	\aB{
		Instead of killing Saul, David steals a little piece of his clothes while he's pooping.
	}
}

\Verse{6}
{
	\hX{וַיְהִי אַחֲרֵי־כֵן וַיַּךְ לֵב־דָּוִד אֹתוֹ עַל אֲשֶׁר כָּרַת אֶת־כָּנָף אֲשֶׁר לְשָׁאוּל׃}
}
{
	\eX{
		But afterward David reproached himself for cutting off the corner
		of Saul’s cloak.
	}
}
{
%	\footnote{
%		Footnote 24:6. YHWH's anointed. David has Saul in his hand and refuses the hand.
%		The phrase protects a king who has become dangerous, which is why the story is
%		more interesting than revenge.
%	}

	\aB{
		Then David feels like he's been a bad David.
	}
}

\Verse{7}
{
	\hX{וַיֹּאמֶר לַאֲנָשָׁיו חָלִילָה לִּי מֵיְהֹוָה אִם־אֶעֱשֶׂה אֶת־הַדָּבָר הַזֶּה לַאדֹנִי לִמְשִׁיחַ יְהֹוָה לִשְׁלֹחַ יָדִי בּוֹ כִּי־מְשִׁיחַ יְהֹוָה הוּא׃}
}
{
	\eX{
		He said to his men, “YHWH forbid that I should do such a thing to
		my lord—YHWH’s anointed—that I should raise my hand against him;
		for he is YHWH’s anointed.”
	}
}
{
	\aB{
		He talks about how it's wrong to steal the king's clothes while he's pooping.
	}

	\aB{
		Convenient opinion for a guy who's about to be king.
	}
}

\Verse{8}
{
	\hX{וַיְשַׁסַּע דָּוִד אֶת־אֲנָשָׁיו בַּדְּבָרִים וְלֹא נְתָנָם לָקוּם אֶל־שָׁאוּל וְשָׁאוּל קָם מֵהַמְּעָרָה וַיֵּלֶךְ בַּדָּרֶךְ׃}
}
{
	\eX{
		David rebuked his men and did not permit them to attack Saul. Saul
		left the cave and started on his way.
	}
}
{
	\aB{
		David scolds his men for being bad.
	}

	\aB{
		Saul finishes pooping and leaves the cave.
	}
}

\Verse{9}
{
	\hX{
		וַיָּקם דָּוִד אַחֲרֵי־כֵן וַיֵּצֵא (מן המערה) [מֵהַמְּעָרָה] וַיִּקְרָא אַחֲרֵי־שָׁאוּל לֵאמֹר אֲדֹנִי הַמֶּלֶךְ וַיַּבֵּט שָׁאוּל אַחֲרָיו וַיִּקֹּד דָּוִד אַפַּיִם אַרְצָה
		וַיִּשְׁתָּחוּ׃
	}
}
{
	\eX{
		Then David also went out of the cave and called after Saul, “My
		lord king!” Saul looked around and David bowed low in homage, with
		his face to the ground.
	}
}
{
	Qere Ketiv.

	\aB{
		David runs after Saul and says ``My lord king!''
	}

	\aB{
		Saul looks around like ``Where'd you come from?''
	}

	\aB{
		David bows.
	}
}

\Verse{10}
{
	\hX{וַיֹּאמֶר דָּוִד לְשָׁאוּל לָמָּה תִשְׁמַע אֶת־דִּבְרֵי אָדָם לֵאמֹר הִנֵּה דָוִד מְבַקֵּשׁ רָעָתֶךָ׃}
}
{
	\eX{
		And David said to Saul, “Why do you listen to the people who say,
		‘David is out to do you harm?’
	}
}
{
	\aB{
		David's like ``Why do you listen to all those liars who say I'm trying to hurt you?''
	}
}

\Verse{11}
{
	\hX{
		הִנֵּה הַיּוֹם הַזֶּה רָאוּ עֵינֶיךָ אֵת אֲשֶׁר־נְתָנְךָ יְהֹוָה הַיּוֹם בְּיָדִי בַּמְּעָרָה וְאָמַר לַהֲרָגְךָ וַתָּחס עָלֶיךָ וָאֹמַר לֹא־אֶשְׁלַח יָדִי בַּאדֹנִי כִּי־מְשִׁיחַ
		יְהֹוָה הוּא׃
	}
}
{
	\eX{
		You can see for yourself now that YHWH delivered you into my hands
		in the cave today. And though I was urged to kill you, I showed
		you pity; for I said, ‘I will not raise a hand against my lord,
		since he is YHWH’s anointed.’
	}
}
{
	\aB{
		I could have killed you in that cave, they told me to, but I said ``No, that's wrong.''
	}
}

\Verse{12}
{
	\hX{
		וְאָבִי רְאֵה גַּם רְאֵה אֶת־כְּנַף מְעִילְךָ בְּיָדִי כִּי בְּכרְתִי אֶת־כְּנַף מְעִילְךָ וְלֹא הֲרַגְתִּיךָ דַּע וּרְאֵה כִּי אֵין בְּיָדִי רָעָה וָפֶשַׁע וְלֹא־חָטָאתִי לָךְ
		וְאַתָּה צֹדֶה אֶת־נַפְשִׁי לְקַחְתָּהּ׃
	}
}
{
	\eX{
		Please, sir, take a close look at the corner of your cloak in my
		hand; for when I cut off the corner of your cloak, I did not kill
		you. You must see plainly that I have done nothing evil or
		rebellious, and I have never wronged you. Yet you are bent on
		taking my life.
	}
}
{
	\aB{
		Look, I stole some of your clothes while you were pooping as proof of how not evil I am.
	}

	\aB{
		But you just keep trying to kill me.
	}
}

\Verse{13}
{
	\hX{יִשְׁפֹּט יְהֹוָה בֵּינִי וּבֵינֶךָ וּנְקָמַנִי יְהֹוָה מִמֶּךָּ וְיָדִי לֹא תִהְיֶה־בָּךְ׃}
}
{
	\eX{
		May YHWH judge between you and me! And may He take vengeance upon
		you for me, but my hand will never touch you.
	}
}
{
	\aB{
		I'll let YHWH get revenge for me against you if he wants.
	}

	\aB{
		Just saying I won't do that.
	}

	\aB{
		But keep that in mind if anything happens.
	}
}

\Verse{14}
{
	\hX{כַּאֲשֶׁר יֹאמַר מְשַׁל הַקַּדְמֹנִי מֵרְשָׁעִים יֵצֵא רֶשַׁע וְיָדִי לֹא תִהְיֶה־בָּךְ׃}
}
{
	\eX{
		As the ancient proverb has it: ‘Wicked deeds come from wicked
		men!’ My hand will never touch you.
	}
}
{
	\aB{
		Theorem 1:
	}

	\aB{
		I'm not a bad guy.
	}

	\aB{
		Proof:
	}

	\aB{
		By the axiom of the ancients:
	}

	\aB{
		Wicked man implies wicked stuff.
	}

	\aB{
		I have done no wicked stuff.
	}

	\aB{
		So by the law of modus tollens, we can conclude that I'm not a wicked man.
	}

	\aB{
		This completes the proof.
	}

	\aB{
		QED.
	}
}

\Verse{15}
{
	\hX{אַחֲרֵי מִי יָצָא מֶלֶךְ יִשְׂרָאֵל אַחֲרֵי מִי אַתָּה רֹדֵף אַחֲרֵי כֶּלֶב מֵת אַחֲרֵי פַּרְעֹשׁ אֶחָד׃}
}
{
	\eX{
		Against whom has the king of Israel come out? Whom are you
		pursuing? A dead dog? A single flea?
	}
}
{
	\aB{
		Theorem 2:
	}

	\aB{
		You're behaving ridiculously.
	}

	\aB{
		Proof:
	}

	\aB{
		Switching now from formal logic to evocative imagery:
	}

	\aB{
		1. Let us take it as given that I'm a very unthreatening animal metaphor. 2. Now, in the category Beh of behaviors, the silliness of a behavior is preserved under maps between a metaphor's source object or domain, and its target object or codomain. 3. Therefore, the diagram commutes, and the theorem follows.
	}

	\aB{
		In summary, pursuing me would be as silly as pursuing those very unthreatening animals, because by (1) I'm like one, and by (2) that's the same.
	}

	\aB{
		This completes the second proof.
	}

	\aB{
		QED.
	}
}

\Verse{16}
{
	\hX{וְהָיָה יְהֹוָה לְדַיָּן וְשָׁפַט בֵּינִי וּבֵינֶךָ וְיֵרֶא וְיָרֵב אֶת־רִיבִי וְיִשְׁפְּטֵנִי מִיָּדֶךָ׃}
}
{
	\eX{
		May YHWH be arbiter and may He judge between you and me! May He
		take note and uphold my cause, and vindicate me against you.”
	}
}
{
	\aB{
		So just to reiterate,
	}

	\aB{
		I'll let YHWH get revenge for me against you if he wants.
	}

	\aB{
		Just saying I won't do that.
	}

	\aB{
		But keep that in mind if anything happens.
	}
}

\Verse{17}
{
	\hX{וַיְהִי כְּכַלּוֹת דָּוִד לְדַבֵּר אֶת־הַדְּבָרִים הָאֵלֶּה אֶל־שָׁאוּל וַיֹּאמֶר שָׁאוּל הֲקֹלְךָ זֶה בְּנִי דָוִד וַיִּשָּׂא שָׁאוּל קֹלוֹ וַיֵּבְךְּ׃}
}
{
	\eX{
		When David finished saying these things to Saul, Saul said, “Is
		that your voice, my son David?” And Saul broke down and wept.
	}
}
{
	\aB{
		So after David says all that, Saul's like:
	}

	\aB{
		``Is that your voice, David?''
	}

	\aB{
		And David's like ``Were you even listening?''
	}

	\aB{
		Then Saul breaks down crying.
	}

	\aB{
		Give him a break, he's getting old.
	}
}

\Verse{18}
{
	\hX{וַיֹּאמֶר אֶל־דָּוִד צַדִּיק אַתָּה מִמֶּנִּי כִּי אַתָּה גְּמַלְתַּנִי הַטּוֹבָה וַאֲנִי גְּמַלְתִּיךָ הָרָעָה׃}
}
{
	\eX{
		He said to David, “You are right, not I; for you have treated me
		generously, but I have treated you badly.
	}
}
{
	\aB{
		Saul agrees with the proofs.
	}

	\aB{
		Saul adds the following corollaries:
	}

	\aB{
		1. David is right. 2. Saul is wrong. 3. David has treated Saul generously. 4. Saul has treated David badly.
	}

	\aB{
		The proofs are left as exercises for the reader.
	}
}

\Verse{19}
{
	\hX{(ואת) [וְאַתָּה] הִגַּדְתָּ הַיּוֹם אֵת אֲשֶׁר־עָשִׂיתָה אִתִּי טוֹבָה אֵת אֲשֶׁר סִגְּרַנִי יְהֹוָה בְּיָדְךָ וְלֹא הֲרַגְתָּנִי׃}
}
{
	\eX{
		Yes, you have just revealed how generously you treated me, for
		YHWH delivered me into your hands and you did not kill me.
	}
}
{
	Qere Ketiv.

	\aB{
		He reiterates that he agrees because of what happened in the pooping scene.
	}
}

\Verse{20}
{
	\hX{וְכִי־יִמְצָא אִישׁ אֶת־אֹיְבוֹ וְשִׁלְּחוֹ בְּדֶרֶךְ טוֹבָה וַיהֹוָה יְשַׁלֶּמְךָ טוֹבָה תַּחַת הַיּוֹם הַזֶּה אֲשֶׁר עָשִׂיתָה לִי׃}
}
{
	\eX{
		If a man meets his enemy, does he let him go his way unharmed?
		Surely, YHWH will reward you generously for what you have done for
		me this day.
	}
}
{
	\aB{
		Theorem 3:
	}

	\aB{
		David does not see Saul as his enemy.
	}

	\aB{
		Proof:
	}

	\aB{
		The proof follows by an application of the rhetorical question strategy. We sketch the proof here, omitting details.
	}

	\aB{
		Assume by way of contradiction that a man meets his enemy and lets him get away unharmed. This contradicts the definition of enemy and completes the proof.
	}

	\aB{
		QED.
	}
}

\Verse{21}
{
	\hX{וְעַתָּה הִנֵּה יָדַעְתִּי כִּי מָלֹךְ תִּמְלוֹךְ וְקָמָה בְּיָדְךָ מַמְלֶכֶת יִשְׂרָאֵל׃}
}
{
	\eX{
		I know now that you will become king, and that the kingship over
		Israel will remain in your hands.
	}
}
{
	\aB{
		Conjecture:
	}

	\aB{
		David will become king, and stay king.
	}

	\aB{
		While the conjecture remains unproved, leading figures in the field of Israelite kingship have expressed confidence in the truth of the proposition.
	}

	\aB{
		(Saul, King. ca. 1000 B.C., personal communication.)
	}
}

\Verse{22}
{
	\hX{וְעַתָּה הִשָּׁבְעָה לִּי בַּיהֹוָה אִם־תַּכְרִית אֶת־זַרְעִי אַחֲרָי וְאִם־תַּשְׁמִיד אֶת־שְׁמִי מִבֵּית אָבִי׃}
}
{
	\eX{
		So swear to me by YHWH that you will not destroy my descendants or
		wipe out my name from my father’s house.”
	}
}
{
	\aB{
		Saul requests that David not kill him and his entire family when he's in charge.
	}
}

\Verse{23}
{
	\hX{וַיִּשָּׁבַע דָּוִד לְשָׁאוּל וַיֵּלֶךְ שָׁאוּל אֶל־בֵּיתוֹ וְדָוִד וַאֲנָשָׁיו עָלוּ עַל־הַמְּצוּדָה׃}
}
{
	\eX{
		David swore to Saul, Saul went home, and David and his men went up
		to the strongholds.
	}
}
{
	\aB{
		David promises.
	}

	\aB{
		He won't.
	}
}
